%%%%%%%%%%%%%%%%%%%%%%%%%%%%%%%%%%%%%%%%%%%%%%%%%%%%%%%%%%%%%%%%%%%%%%%%%%%%
% AGUJournalTemplate.tex: this template file is for articles formatted with LaTeX
%
% This file includes commands and instructions
% given in the order necessary to produce a final output that will
% satisfy AGU requirements, including customized APA reference formatting.
%
% You may copy this file and give it your
% article name, and enter your text.
%
% guidelines and troubleshooting are here: 

%% To submit your paper:
\documentclass[draft]{agujournal2019}
\usepackage{url} %this package should fix any errors with URLs in refs.
\usepackage{lineno}
\usepackage[inline]{trackchanges} %for better track changes. finalnew option will compile document with changes incorporated.
\usepackage{soul}
\linenumbers
\graphicspath{{figures/}} % path to folder with figures
\usepackage{float}
%%%%%%%
% As of 2018 we recommend use of the TrackChanges package to mark revisions.
% The trackchanges package adds five new LaTeX commands:
%
%  \note[editor]{The note}
%  \annote[editor]{Text to annotate}{The note}
%  \add[editor]{Text to add}
%  \remove[editor]{Text to remove}
%  \change[editor]{Text to remove}{Text to add}
%
% complete documentation is here: http://trackchanges.sourceforge.net/
%%%%%%%

\draftfalse

\journalname{Journal of Advances in Modeling Earth Systems (JAMES)}

\begin{document}

%%%%%%%%%%%%%%%%%%%%%%%%%%%%%%%%%%%%%%%%%%%%%%%

\title{Ensemble Learning Methods for Image Classification Using Multitude Decision Trees in GRASS GIS for Mapping Seasonal Fluctuations of Etosha Pan, Namibia}

%%%%%%%%%%%%%%%%%%%%%%%%%%%%%%%%%%%%%%%%%%%%%%%

\authors{Polina Lemenkova\affil{1}\thanks{Schillerstraße 30, A-5020 Salzburg, Austria}}

\affiliation{1}{Department of Geoinformatics, Faculty of Digital and Analytical Sciences, Universität Salzburg. Schillerstraße 30, A-5020 Salzburg, Austria}

\correspondingauthor{Polina Lemenkova}{polina.lemenkova@plus.ac.at}

%%%%%%%%%%%%%%%%%%%%%%%%%%%%%%%%%%%%%%%%%%%%%%%

\begin{keypoints}
\item Random Forest
\item Remote Sensing
\item Image Processing
\end{keypoints}

%%%%%%%%%%%%%%%%%%%%%%%%%%%%%%%%%%%%%%%%%%%%%%%
%
%  ABSTRACT and PLAIN LANGUAGE SUMMARY
%
% A good Abstract will begin with a short description of the problem
% being addressed, briefly describe the new data or analyses, then
% briefly states the main conclusion(s) and how they are supported and
% uncertainties.

% The Plain Language Summary should be written for a broad audience,
% including journalists and the science-interested public, that will not have 
% a background in your field.
%
% A Plain Language Summary is required in GRL, JGR: Planets, JGR: Biogeosciences,
% JGR: Oceans, G-Cubed, Reviews of Geophysics, and JAMES.
% see http://sharingscience.agu.org/creating-plain-language-summary/)
%
%%%%%%%%%%%%%%%%%%%%%%%%%%%%%%%%%%%%%%%%%%%%%%%

%% \begin{abstract} starts the second page

\begin{abstract}
Identification of fluctuations in Etosha salt pan poses a challenge in the environmental domain. Etosha Pan presents unique salt landscape and noticeable geological feature in northern Namibia. With seasonal fluctuation cycles from dry to wet periods, the water level in the basin contrasts from evaporated to flooding. Such changes affect vegetation and surrounding landscapes. Monitoring dynamics of the Etosha salt pan is possible using a time series analysis based on the Remote Sensing (RS) data. Image classification is the most common type of RS data processing which can be used for cartographic visualization of changes in land cover types. Recently, the methods of Machine Learning (ML) experience significant development which paved the way to their effective applications in Earth system data processing. Ensemble learning approach with Random Forest (RF) presents the example of such advanced methods applied to satellite image classification. It is based on the integration of several ML approaches including the decision trees algorithm with controlled variance which improves the accuracy of image partition process. This paper improves the image classification methods using training algorithm of RF by the modules of GRASS GIS based on the embedded Scikit-Learn machine learning library of Python. Several existing techniques of GRASS GIS were investigated and tested, with a hybrid approach using a combination of RS modules such as 'r.learn.predict', 'r.random', 'r.learn.train', 'r.category', 'i.group', among others. As a result, land cover types around Etosha Pan were identified using six multi-spectral Landsat 8-9 OLI/TIRS images. The classification of images collected in different seasons for consecutive months from March to August enabled to monitor the process of slow evaporation in Etosha Pan during the transitional wet-to-dry period. The comparison of images revealed the distinct variations in the surface of basin from wet period in March to the dry period when the basin is covered by the crust of salt and minerals. The series of the presented maps made using ML-based methods of GRASS GIS demonstrates the opportunities of the ensemble learning tools to image processing by RF classifiers. Such novel approaches enable accurate classification of geospatial data for environmental analysis and interpretation of landscape dynamics. The proposed method of ML-based image classification in GRASS GIS and its efficacy in RS data evaluation proposes an accurate tool for forecasting and prognosis of changes in south African landscapes of Namibia.
\end{abstract}

\section*{Plain Language Summary}
Etosha pan located is a large salt pan in northern Namibia. It is a vast basin in the land surface where water collects seasonally depending on rains and evaporates during dry periods with the surface covered by salt and minerals. Satellite image classification presents effective method for monitoring these changes using a comparison of several scenes collected during dry and wet periods. While traditional classification methods is usually based on the Geographic Information Systems (GIS) tools, a novel tools of Machine Learning (ML) presents a more effective, automated and accurate approach to image processing. This paper presents the application of one of the ML methods using Random Forest (RF) algorithms of GRASS GIS which utilises the embedded libraries of Python Scikit-Learn for image processing. The results show landscape dynamics in Etosha Pan from March to August in 2022 using a series of Landsat satellite images classified by RF method.

%%%%%%%%%%%%%%%%%%%%%%%%%%%%%%%%%%%%%%%%%%%%%%%
%
%  BODY TEXT
%
%%%%%%%%%%%%%%%%%%%%%%%%%%%%%%%%%%%%%%%%%%%%%%%

%%% Suggested section heads:
 \section{Introduction}
%
% The main text should start with an introduction. Except for short
% manuscripts (such as comments and replies), the text should be divided
% into sections, each with its own heading.

% Headings should be sentence fragments and do not begin with a
% lowercase letter or number. Examples of good headings are:

 \section{Materials and Methods}
\subsection{Study Area}
The study area includes the northern part of Namibia, the Etosha Pan, Figure \ref{fig01}.
 
 \begin{figure}[H]
	\noindent\includegraphics[width=\textwidth]{Fig_01.jpg}
\caption{Topographic map of Namibia with indicated study area: Etosha Pan. Mapping software: Generic Mapping Tools (GMT). Map source: author.}
\label{fig01}
\end{figure}

 \begin{figure}[H]
	\noindent\includegraphics[width=\textwidth]{Fig_02.jpg}
\caption{Land cover types of Namibia according to FAO. Mapping software: QGIS. Map source: author.}
\label{fig02}
\end{figure}

 \subsection{Random Forest (RF)}

A Random Forest (RF) function algorithm of ensemble learning was derived from GRASS GIS and used to test the simulation of the satellite image time series on Etosha Pan region, Namibia. The trained pixels were obtained from the 'RandomForestClassifier' algorithms connected to the Scikit-Learn libraries of Python embedded with GRASS GIS ML function of image analysis.  The RF is a type of bagging algorithm of image classification that uses the Machine Learning (ML) methods of the classification and regression tree decision as learners. The randomness of data points in a raster image evaluated in this algorithm is assessed for reducing model variance. As a result, the advantages of the RF consists in excellent generalisation approach to data classification which combines the results of multiple decision trees to reach a single classification result through anti-overfitting.

\section{Data}
The data used in this study includes four satellite images of the Landsat 8-9 Operational Land Imager and Thermal Infrared Sensor (Landsat 8-9 OLI/TIRS). The original images were processed and stored in Universal Transverse Mercator (UTM) Zone 33 for northern Namibia. The dataset is presented in Figure \ref{fig03}.

\begin{figure}[H]
  \centering
  \begin{tabular}[b]{c}
    \includegraphics[width=4cm]{Fig_03a_mar.jpg} \\
    \small (a)
  \end{tabular}
  \begin{tabular}[b]{c}
    \includegraphics[width=4cm]{Fig_03b_apr.jpg} \\
    \small (b)
  \end{tabular}
  \begin{tabular}[b]{c}
    \includegraphics[width=4cm]{Fig_03c_may.jpg} \\
    \small (c)
     \end{tabular}
     \begin{tabular}[b]{c}
    \includegraphics[width=4cm]{Fig_03d_jun.jpg} \\
    \small (d)
  \end{tabular}
  \begin{tabular}[b]{c}
    \includegraphics[width=4cm]{Fig_03e_jul.jpg} \\
    \small (e)
  \end{tabular}
  \begin{tabular}[b]{c}
    \includegraphics[width=4cm]{Fig_03f_aug.jpg} \\
    \small (f)
  \end{tabular}
  \caption{Landsat 8-9 OLI/TIRS images used as dataset covering the Etosha Pan region, Namibia: \textbf{(a)}: 26 March 2022;\textbf{(b)}: 11 April 2022; \textbf{(c)}: 5 May 2022; \textbf{(d)}: 6 June 2022; \textbf{(e)}: 16 July 2022; \textbf{(f)}: 9 August 2022. The images are shown in RGB natural colours. Data sources: USGS}
\end{figure}

\begin{figure}[H]
  \centering
  \begin{tabular}[b]{c}
    \includegraphics[width=4cm]{Fig_04a_754.jpg} \\
    \small (a)
  \end{tabular}
  \begin{tabular}[b]{c}
    \includegraphics[width=4cm]{Fig_04b_234.jpg} \\
    \small (b)
  \end{tabular}
  \begin{tabular}[b]{c}
    \includegraphics[width=4cm]{Fig_04c_345.jpg} \\
    \small (c)
     \end{tabular}
    \begin{tabular}[b]{c}
    \includegraphics[width=4cm]{Fig_04d_753.jpg} \\
    \small (d)
  \end{tabular}
  \begin{tabular}[b]{c}
    \includegraphics[width=4cm]{Fig_04e_345.jpg} \\
    \small (e)
  \end{tabular}
  \begin{tabular}[b]{c}
    \includegraphics[width=4cm]{Fig_04f_234.jpg} \\
    \small (f)
     \end{tabular}
     \begin{tabular}[b]{c}
    \includegraphics[width=4cm]{Fig_04g_357.jpg} \\
    \small (g)
  \end{tabular}
  \begin{tabular}[b]{c}
    \includegraphics[width=4cm]{Fig_04h_156} \\
    \small (h)
  \end{tabular}
  \begin{tabular}[b]{c}
    \includegraphics[width=4cm]{Fig_04i_741.jpg} \\
    \small (i)
     \end{tabular}
  \caption{Examples of different combinations of the created colour composites highlighting the Etosha Pan, Namibia. The images are generated based on triplets of Landsat 8-9 OLI/TIRS bands on 2022, using 'r.composite' module of GRASS GIS: \textbf{(a)}: March 2022: Bands 7-5-3;\textbf{(b)}: March 2022: Bands 2-3-4; \textbf{(c)}: March 2022: Bands 3-4-5; \textbf{(d)}: April 2022: Bands 7-5-4;\textbf{(e)}: April 2022: Bands 3-4-5; \textbf{(f)}: April 2022: Bands 3-4-5; \textbf{(g)}: May 2022: Bands 3-5-7; \textbf{(h)}: June 2022: Bands 1-5-6; \textbf{(i)}: August 2022: Bands 7-4-1.}
\end{figure}

\section{Results}

\begin{figure}[H]
  \centering
  \begin{tabular}[b]{c}
    \includegraphics[width=4.1cm]{Fig_05a.jpg} \\
    \small (a)
  \end{tabular}
  \begin{tabular}[b]{c}
    \includegraphics[width=4.1cm]{Fig_05b.jpg} \\
    \small (b)
  \end{tabular}
  \begin{tabular}[b]{c}
    \includegraphics[width=4.1cm]{Fig_05c.jpg} \\
    \small (c)
     \end{tabular}
     \begin{tabular}[b]{c}
    \includegraphics[width=4.1cm]{Fig_05d.jpg} \\
    \small (d)
  \end{tabular}
  \begin{tabular}[b]{c}
    \includegraphics[width=4.1cm]{Fig_05e.jpg} \\
    \small (e)
  \end{tabular}
  \begin{tabular}[b]{c}
    \includegraphics[width=4.1cm]{Fig_05f.jpg} \\
    \small (f)
  \end{tabular}
  \caption{Classification of the Landsat-8-9 OLI/TIRS images covering the Etosha Pan region, Namibia using 'MaxLike' method: \textbf{(a)}: 26 March 2022;\textbf{(b)}: 11 April 2022; \textbf{(c)}: 5 May 2022; \textbf{(d)}: 6 June 2022; \textbf{(e)}: 16 July 2022; \textbf{(f)}: 9 August 2022.}
\end{figure}

\begin{figure}[H]
  \centering
  \begin{tabular}[b]{c}
    \includegraphics[width=4.1cm]{Fig_06a.jpg} \\
    \small (a)
  \end{tabular}
  \begin{tabular}[b]{c}
    \includegraphics[width=4.1cm]{Fig_06b.jpg} \\
    \small (b)
  \end{tabular}
  \begin{tabular}[b]{c}
    \includegraphics[width=4.1cm]{Fig_06c.jpg} \\
    \small (c)
     \end{tabular}
     \begin{tabular}[b]{c}
    \includegraphics[width=4.1cm]{Fig_06d.jpg} \\
    \small (d)
  \end{tabular}
  \begin{tabular}[b]{c}
    \includegraphics[width=4.1cm]{Fig_06e.jpg} \\
    \small (e)
  \end{tabular}
  \begin{tabular}[b]{c}
    \includegraphics[width=4.1cm]{Fig_06f.jpg} \\
    \small (f)
  \end{tabular}
  \caption{Accuracy assessment using the rejection probability for pixels classified for Landsat-8-9 OLI/TIRS images on the Etosha Pan region, Namibia: \textbf{(a)}: 26 March 2022;\textbf{(b)}: 11 April 2022; \textbf{(c)}: 5 May 2022; \textbf{(d)}: 6 June 2022; \textbf{(e)}: 16 July 2022; \textbf{(f)}: 9 August 2022.}
\end{figure}

\begin{figure}[H]
  \centering
  \begin{tabular}[b]{c}
    \includegraphics[width=4.1cm]{Fig_07a.jpg} \\
    \small (a)
  \end{tabular}
  \begin{tabular}[b]{c}
    \includegraphics[width=4.1cm]{Fig_07b.jpg} \\
    \small (b)
  \end{tabular}
  \begin{tabular}[b]{c}
    \includegraphics[width=4.1cm]{Fig_07c.jpg} \\
    \small (c)
     \end{tabular}
     \begin{tabular}[b]{c}
    \includegraphics[width=4.1cm]{Fig_07d.jpg} \\
    \small (d)
  \end{tabular}
  \begin{tabular}[b]{c}
    \includegraphics[width=4.1cm]{Fig_07e.jpg} \\
    \small (e)
  \end{tabular}
  \begin{tabular}[b]{c}
    \includegraphics[width=4.1cm]{Fig_07f.jpg} \\
    \small (f)
  \end{tabular}
  \caption{Results of the Random Forest (RF) classification using ML approach of GRASS GIS for the Landsat-8-9 OLI/TIRS images on Etosha Pan, Namibia: \textbf{(a)}: 26 March 2022;\textbf{(b)}: 11 April 2022; \textbf{(c)}: 5 May 2022; \textbf{(d)}: 6 June 2022; \textbf{(e)}: 16 July 2022; \textbf{(f)}: 9 August 2022.}
\end{figure}

\section{Conclusions}

A novel hybrid ensemble learning ML method of GRASS GIS to image classification is presented in this manuscript. The approach was employed to classify a series of multispectral Landsat 8-9 OLI/TIRS images in order to indicate, detect and visualise the dynamics of water evaporation and the increase in salt pan in the Etosha basin, as well as to demonstrate changes in the surrounding lacustrine landscapes. In this paper, there are major three steps of image processing: clustering with generating training pixels, traditional classification using maximum Likelihood approach and Random Forest (RF) classification of ML. Additionally, the colour composites were created using several bands composites of Landsat. The classification accuracy was assessed using the computed confidence level of pixels assigned into ten major land cover classes around Etosha Pan. The RF algorithm provides the output of multiple decision trees to reach a classification of the satellite image. Thus, the hybrid classifier of prediction pattern evaluates the assignment of pixels on a raster scene using multiple decision trees.

The efficiency of environmental mapping that operate with the satellite images, especially for monitoring seasonal fluctuations in land cover types can now be improved by using the Machine Learning (ML) methods of Python integrated with cartographic tools. Thus, the analysis of scenes indicates that the evaporation of water during dry periods coincides with the cycle of flooding and evaporation in the north Namibia. During dry period, a mineral-encrusted salt pan covers the surface of basin which reaches its peak in August while the most water-filled period is in March. The comparison of these scenes as ML-classified satellite images show gradual development in salinization of the Etosha Pan basin during six consecutive months in 2022. The results of experiments demonstrate that RF has a high automation in image processing, accuracy in pixel assignments, increased details of classification with more fragmented identified patterns in the landscapes of norther Namibia. 

The RF method of image processing is pitted against the other traditional classification methods such as k-means Clustering and Maximum Likelihood (MaxLike) of GRASS GIS. The use of such advanced solution helps in providing automated tools of satellite image processing and classification. Moreover, the feasibility of the RF algorithms enables to use it as an innovative approach that may be used in practise for diverse types of satellite images, including Sentinel, SPOT and other imagery. The outcome of the framework in the GRASS GIS based image processing using ensemble learning is the output of the series of the classified satellite images processed using decision criteria analysis of ML algorithms. To enhance its application in similar future studies, the RF classifiers technique can be applies as effective ML tools of the GRASS GIS modules to processing Earth observation data for operative environmental monitoring.

%
% ---------------
% EXAMPLE TABLE
%
% \begin{table}
% \caption{Time of the Transition Between Phase 1 and Phase 2$^{a}$}
% \centering
% \begin{tabular}{l c}
% \hline
%  Run  & Time (min)  \\
% \hline
%   $l1$  & 260   \\
%   $l2$  & 300   \\
%   $l3$  & 340   \\
%   $h1$  & 270   \\
%   $h2$  & 250   \\
%   $h3$  & 380   \\
%   $r1$  & 370   \\
%   $r2$  & 390   \\
% \hline
% \multicolumn{2}{l}{$^{a}$Footnote text here.}
% \end{tabular}
% \end{table}

%%%%%%%%%%%%%%%%%%%%%%%%%%%%%%%%%%%%%%%%%%%%%%%
% SIDEWAYS FIGURES and TABLES
% AGU prefers the use of {sidewaystable} over {landscapetable} as it causes fewer problems.
%
% \begin{sidewaysfigure}
% \includegraphics[width=20pc]{figsamp}
% \caption{caption here}
% \label{newfig}
% \end{sidewaysfigure}
%
%  \begin{sidewaystable}
%  \caption{Caption here}
% \label{tab:signif_gap_clos}
%  \begin{tabular}{ccc}
% one&two&three\\
% four&five&six
%  \end{tabular}
%  \end{sidewaystable}

%%% End of body of article

%%%%%%%%%%%%%%%%%%%%%%%%%%%%%%%%%%%%%%%%%%%%%%%
%% Optional Appendices go here
%
% The \appendix command resets counters and redefines section heads
%
% After typing \appendix
%
%\section{Here Is Appendix Title}
% will show
% A: Here Is Appendix Title
%
%\appendix
%\section{Here is a sample appendix}


%%%% NOTE that acronyms in the final published version will be spelled out when used in figure captions.
 \begin{acronyms}
   \acro{USGS}
    United States Geological Survey\\
   \acro{GMT}
    Generic Mapping Tools \\
   \acro{Landsat OLI/TIRS }
   Landsat Operational Land Imager and Thermal Infrared Sensor \\
   \acro{FAO}
   Food and Agriculture Organization \\
\end{acronyms}


%%%%%%%%%%%%%%%%%%%%%%%%%%%%%%%%%%%%%%%%%%%%%%%
% Notation
%   \begin{notation}
%   \notation{$a+b$} Notation Definition here
%   \notation{$e=mc^2$}
%   Equation in German-born physicist Albert Einstein's theory of special
%  relativity that showed that the increased relativistic mass ($m$) of a
%  body comes from the energy of motion of the body—that is, its kinetic
%  energy ($E$)—divided by the speed of light squared ($c^2$).
%   \end{notation}


%%%%%%%%%%%%%%%%%%%%%%%%%%%%%%%%%%%%%%%%%%%%%%%
%
% DATA SECTION and ACKNOWLEDGMENTS
%
%%%%%%%%%%%%%%%%%%%%%%%%%%%%%%%%%%%%%%%%%%%%%%%

\section*{Open Research Section}
This section MUST contain a statement that describes where the data supporting the conclusions can be obtained. Data cannot be listed as ''Available from authors'' or stored solely in supporting information. Citations to archived data should be included in your reference list. Wiley will publish it as a separate section on the paper’s page. Examples and complete information are here:
https://www.agu.org/Publish with AGU/Publish/Author Resources/Data for Authors


%%%%%%%%%%%%%%%%%%%%%%%%%%%%%%%%%%%%%%%%%%%%%%%
% REFERENCES and BIBLIOGRAPHY
%
\bibliography{Namibia}
\nocite{*}
% don't specify bibliographystyle
%
%%%%%%%%%%%%%%%%%%%%%%%%%%%%%%%%%%%%%%%%%%%%%%%

%\bibliography{ enter your bibtex bibliography filename here }



%Reference citation instructions and examples:
%
% Please use ONLY \cite and \citeA for reference citations.
% \cite for parenthetical references
% ...as shown in recent studies (Simpson et al., 2019)
% \citeA for in-text citations
% ...Simpson et al. (2019) have shown...
%
%
%...as shown by \citeA{jskilby}.
%...as shown by \citeA{lewin76}, \citeA{carson86}, \citeA{bartoldy02}, and \citeA{rinaldi03}.
%...has been shown \cite{jskilbye}.
%...has been shown \cite{lewin76,carson86,bartoldy02,rinaldi03}.
%... \cite <i.e.>[]{lewin76,carson86,bartoldy02,rinaldi03}.
%...has been shown by \cite <e.g.,>[and others]{lewin76}.
%
% apacite uses < > for prenotes and [ ] for postnotes
% DO NOT use other cite commands (e.g., \citet, \citep, \citeyear, \nocite, \citealp, etc.).
%

\end{document}



More Information and Advice:

%%%%%%%%%%%%%%%%%%%%%%%%%%%%%%%%%%%%%%%%%%%%%%%
%
%  SECTION HEADS
%
%%%%%%%%%%%%%%%%%%%%%%%%%%%%%%%%%%%%%%%%%%%%%%%

% Capitalize the first letter of each word (except for
% prepositions, conjunctions, and articles that are
% three or fewer letters).

% AGU follows standard outline style; therefore, there cannot be a section 1 without
% a section 2, or a section 2.3.1 without a section 2.3.2.
% Please make sure your section numbers are balanced.
% ---------------
% Level 1 head
%
% Use the \section{} command to identify level 1 heads;
% type the appropriate head wording between the curly
% brackets, as shown below.
%
%An example:
%\section{Level 1 Head: Introduction}
%
% ---------------
% Level 2 head
%
% Use the \subsection{} command to identify level 2 heads.
%An example:
%\subsection{Level 2 Head}
%
% ---------------
% Level 3 head
%
% Use the \subsubsection{} command to identify level 3 heads
%An example:
%\subsubsection{Level 3 Head}
%
%---------------
% Level 4 head
%
% Use the \subsubsubsection{} command to identify level 3 heads
% An example:
%\subsubsubsection{Level 4 Head} An example.
%
%%%%%%%%%%%%%%%%%%%%%%%%%%%%%%%%%%%%%%%%%%%%%%%
%
%  IN-TEXT LISTS
%
%%%%%%%%%%%%%%%%%%%%%%%%%%%%%%%%%%%%%%%%%%%%%%%
%
% Do not use bulleted lists; enumerated lists are okay.
% \begin{enumerate}
% \item
% \item
% \item
% \end{enumerate}
%


%%%%%%%%%%%%%%%%%%%%%%%%%%%%%%%%%%%%%%%%%%%%%%%
%
%  EQUATION NUMBERING: COUNTER
%
%%%%%%%%%%%%%%%%%%%%%%%%%%%%%%%%%%%%%%%%%%%%%%%

% You may change equation numbering by resetting
% the equation counter or by explicitly numbering
% an equation.

% To explicitly number an equation, type \eqnum{}
% (with the desired number between the brackets)
% after the \begin{equation} or \begin{eqnarray}
% command.  The \eqnum{} command will affect only
% the equation it appears with; LaTeX will number
% any equations appearing later in the manuscript
% according to the equation counter.
%

% If you have a multiline equation that needs only
% one equation number, use a \nonumber command in
% front of the double backslashes (\\) as shown in
% the multiline equation above.

% If you are using line numbers, remember to surround
% equations with \begin{linenomath*}...\end{linenomath*}

%  To add line numbers to lines in equations:
%  \begin{linenomath*}
%  \begin{equation}
%  \end{equation}
%  \end{linenomath*}



